%% kapa: exam=1 round=37 year=2026 period=1 form=A subject=real_estate subject_ko=부동산학원론 seq=38 number=118
%% source: Q-Net 감정평가사 기출문제 (cst00309) artlSeq=5258924 file=raw/1008/5258924/2026년 제37회 감정평가사 1차 1교시.pdf
%% conversion: kordoc 4.24.51 -> tools/gen1.py (auto); review=proofread-sonnet; flags=none
\begin{question}{118}
감정평가에 관한 규칙상 물건별 평가방법에 관한 내용으로 옳지 않은 것은?
\begin{choices}
  \item 산림을 감정평가할 때에 산지와 입목(立木)을 구분하여 감정평가해야 한다. 이 경우 입목은 거래사례비교법을 적용하되, 소경목림(小徑木林)인 경우에는 원가법을 적용할 수 있다.
  \item 과수원을 감정평가할 때에 거래사례비교법을 적용해야 한다.
  \item 항공기를 감정평가할 때에 원가법을 적용해야 한다.
  \item 광업권을 감정평가할 때에 광업재단의 감정평가액에서 해당 광산의 현존시설 가액을 빼고 감정평가해야 한다. 이 경우 광산의 현존시설 가액은 적정 생산규모와 가행조건(稼行條件) 등을 고려하여 산정하되 과잉유휴시설을 포함하여 산정한다.
  \item 소음ㆍ진동ㆍ일조침해 또는 환경오염 등(이하 “소음등”이라 한다)으로 대상물건의 가치가 하락한 경우 그 가치하락분을 감정평가할 때에 소음등이 발생하기 전의 대상물건의 가액 및 원상회복비용 등을 고려해야 한다.
\end{choices}
\end{question}
\answer{4}
